% !TEX root = chapter-standalone.tex

\section[Product-Hom adjunction]{Example of a ``Product-Hom'' adjunction}
We will consider an \SY{adjunction} between the category \Set and itself which is a basic representative of a certain ``type'' of \SY{adjunction} that appears all over mathematics.
This type of \SY{adjunction} might be called a ``Product-Hom'' \SY{adjunction}.

Fix a set~\setB and consider \SY{functors}~$\funa$ and~$\funb$
\begin{equation}
\begin{tikzcd}[every arrow/.append style={arrowfunctorstyle}]
    {\Set}
     &  &  &  & \Set
    \arrow[""{name=0, anchor=center, inner sep=0}, "\funa", bend left=30, from=1-1, to=1-5]
    \arrow[""{name=1, anchor=center, inner sep=0}, "\funb", bend left=30, from=1-5, to=1-1]
\end{tikzcd}
\end{equation}
which are defined as follows.

On the level of objects -- in this case these are sets -- we define
\begin{equation}\label{eq:product-with-a-set-functor-objects}
    \funa(\setA) =  \setA \cartprod \setB 
\end{equation}
and
\begin{equation}\label{eq:set-expontiation-functor-objects}
    \funb(\setA) =  \setA^\setB
\end{equation}
(recall that  $\setA^\setB= \HomSet{\Set}{\setB}{\setA}$ denotes the set of functions from $\setB$ to $\setA$.)

\todotextjira{474}{\alphubel: @JL: too synthentic, not readable}
Now let's define $\funa$ and $\funb$ on the level of morphisms (which, in this case, are functions). Given a function~$\mora\colon \setA \mto \setAprime$, we define
\begin{equation}\label{eq:product-with-a-set-functor-morphisms}
    \begin{aligned}
        \funa(\mora)\colon \quad \setA \cartprod \setB \quad & \sto \quad \setAprime \cartprod \setB \\
        \tup{\ela, \elb}  \quad          & \mapsto \quad \tup{\mora(\ela), \elb} .
    \end{aligned}
\end{equation}
%\begin{equation}\label{eq:product-with-a-set-functor-morphisms}
%    \funa(\mora) =   \mora \funcprod \catidat\setB
%\end{equation}
and
\begin{equation}\label{eq:set-exponentiation-functor-morphisms}
    \begin{aligned}
        \funb(\mora)\colon \quad \setA^\setB \quad & \sto \quad \setAprime^\setB \\
        \morb \quad                         & \mapsto \quad \morb \mthen \mora.
    \end{aligned}
\end{equation}

A useful shorthand notation for these functors is $\funa = ( - ) \cartprod \setB$ and $\funb = ( - )^\setB$. 
%
They form an adjunction
\begin{equation}
\begin{tikzcd}[every arrow/.append style={arrowfunctorstyle}]
    {\Set}
     &  &  &  & \Set
    \arrow[""{name=0, anchor=center, inner sep=0}, "{( - ) \cartprod \setB}", bend left=30, from=1-1, to=1-5]
    \arrow[""{name=1, anchor=center, inner sep=0}, "{( - )^\setB}", bend left=30, from=1-5, to=1-1]
    \arrow["\dashv"{anchor=center, rotate=-90}, draw=none, from=0, to=1]
\end{tikzcd}
\end{equation}
Let's see how, according to each of the two different definitions of adjunction. 

\subsection{Checking the hom-set definition of adjunction}

Consider \cref{def:adj-iso}. In the present example, there is a natural isomorphism
\begin{equation}\label{eq:prod-hom-adjunction-transpose-nat-iso}
    \stylenat{\tau}\colon \quad \HomSet{\Set}{\funa(- )}{-} \quad \ntolong \quad \HomSet{\Set}{-}{\funb( - )}
\end{equation}
whose component at~$\tup{\setA,\setC}$ is the isomorphism
\begin{equation}\label{eq:prod-hom-adjunction-transpose-components}
    \stylenat{\tau}_{\setA,\setC} \colon \quad  \HomSet\Set{ \setA \cartprod \setB}{\setC} \quad \mto \quad \HomSet{\Set}{\setA}{\setC^\setB}
\end{equation}
given by ``partial evaluation''.
Namely, given~$\mora \colon   \setA \cartprod \setB \mto \setC$, this is mapped by~$\stylenat{\tau}_{\setA,\setC}$ to the function
\begin{equation}
    \begin{aligned}
        \stylenat{\tau}_{\setA,\setC}(\mora) \colon \quad \setA \quad & \to \quad  \setC^\setB \\
        \setAel          \quad             & \mapsto \quad \mora(\setAel, - ).
    \end{aligned}
\end{equation}

\subsection{Checking the (co)unit definition of adjunction}

Now consider the (co)unit definition of adjunctions. In the present example, we claim that unit and co-unit natural transformations 
\begin{equation}
    \funid_\Set \ \overset{\equivunit}{\ntolong}\ \funa \fthen \funb \qqand \funb \fthen \funa \ \overset{\equivcounit}{\ntolong}\ \funid_\Set.
\end{equation}
are defined as follows. In components, for any set \setA, we define
\begin{equation}\label{eq:prod-hom-adjunction-unit}
    \begin{aligned}
        \equivunit_\setA \colon \quad \setA \quad & \to  \quad (  \setA \cartprod \setB)^\setB \\
        \setAel     \quad                  & \mapsto \quad (\setBel \mapsto \tup{\setAel,\setBel})
    \end{aligned}
\end{equation}
and
\begin{equation}\label{eq:prod-hom-adjunction-co-unit}
    \begin{aligned}
        \equivcounit_\setA \colon \quad  (\setA^\setB) \cartprod \setB \quad & \sto \quad \setA \\
        \tup{\mora,\setBel} \quad                                    & \mapsto \quad \mora(\setBel).
    \end{aligned}
\end{equation}


